\documentclass[aspectratio=169]{beamer}
\usetheme{Madrid}
\usecolortheme{default}

\usepackage{listings}
\usepackage{xcolor}
\usepackage{hyperref}


%\lstdefinestyle{sparkstyle}{
%	language=Python,
%	basicstyle=\ttfamily\scriptsize,
%	keywordstyle=\bfseries,
%	commentstyle=\itshape,
%	stringstyle=\ttfamily,
%	showstringspaces=false,
%	breaklines=true,
%	breakatwhitespace=true,
%	frame=single,
%	numbers=left,
%	numberstyle=\tiny,
%	tabsize=4
%}


% Code listing style for Python/Spark
\lstdefinestyle{sparkstyle}{
	language=Python,
	basicstyle=\ttfamily\scriptsize,
	keywordstyle=\color{blue!80!black},
	stringstyle=\color{red!70!black},
	commentstyle=\color{green!50!black},
	numbers=left,
	numberstyle=\tiny\color{gray},
	stepnumber=1,
	numbersep=8pt,
	backgroundcolor=\color{gray!5},
	frame=single,
	rulecolor=\color{gray!40},
	breaklines=true,
	showstringspaces=false,
	tabsize=2,
	captionpos=b
}

\lstdefinestyle{shellstyle}{
	language=bash,
	basicstyle=\ttfamily\scriptsize,
	keywordstyle=\color{blue!80!black},
	stringstyle=\color{red!70!black},
	commentstyle=\color{green!50!black},
	backgroundcolor=\color{gray!5},
	frame=single,
	rulecolor=\color{gray!40},
	breaklines=true,
	showstringspaces=false
}

\title{Fixing the Spark Notebook}
\subtitle{From Timeouts to Tuned Configurations}
\author{Data Engineering Course}
\date{\today}

\begin{document}
	
	% ---------------------------------------------------------------
	% Title Slide
	% ---------------------------------------------------------------
	\begin{frame}
		\titlepage
	\end{frame}
	
	% ---------------------------------------------------------------
	\begin{frame}{Outline}
		\tableofcontents
	\end{frame}
	
	% ---------------------------------------------------------------
	\section{The Problem: Notebook Timeouts}
	% ---------------------------------------------------------------
	
	\begin{frame}{The Problem: Notebook Timeouts}
		\begin{block}{Symptoms}
			\begin{itemize}
				\item Running the Spark notebook \textbf{as-is} takes \textbf{4--6 hours}.
				\item Eventually it \textbf{times out} or loses track.
				\item Not exactly a crash --- it silently hangs after several hours.
				\item Hard to see what is happening at that point.
			\end{itemize}
		\end{block}
		
		\begin{alertblock}{Impact}
			\begin{itemize}
				\item No useful output after a full working day.
				\item Iteration is painfully slow.
				\item We need a more reliable and faster execution.
			\end{itemize}
		\end{alertblock}
	\end{frame}
	
	% ---------------------------------------------------------------
	\begin{frame}{Why Does It Time Out?}
		\begin{itemize}
			\item The notebook processes a \textbf{large amount of data}.
			\item Default Spark configuration may not provide enough:
			\begin{itemize}
				\item Executor memory
				\item Number of executors
				\item Driver memory
				\item CPU cores
			\end{itemize}
			\item When resources are insufficient, Spark spills to disk, retries, and slows down dramatically.
			\item After 5--6 hours, the session may be killed by the cluster manager or simply become unresponsive.
		\end{itemize}
		
		\begin{exampleblock}{Key Insight}
			We need to \textbf{change the configuration} --- executor memory, number of executors, and possibly driver memory --- to match the workload.
		\end{exampleblock}
	\end{frame}
	
	% ---------------------------------------------------------------
	\section{How We Got the Idea: Old Notebooks}
	% ---------------------------------------------------------------
	
	\begin{frame}{Learning from Old Notebooks}
		\begin{itemize}
			\item We examined \textbf{3--5 year old notebooks} in the repository.
			\item Those notebooks were \textbf{manually configuring} Spark:
			\begin{itemize}
				\item \texttt{spark.executor.memory}
				\item \texttt{spark.executor.cores}
				\item \texttt{spark.executor.instances}
				\item \texttt{spark.driver.memory}
			\end{itemize}
			\item They were designed to run \textbf{large workloads} and completed successfully.
			\item This gave us the idea: \textbf{replicate their configuration approach}.
		\end{itemize}
		
		\begin{block}{Observation}
			Older code did not have the luxury of pre-configured ``large'' notebooks.\\
			Engineers explicitly tuned Spark for each job.
		\end{block}
	\end{frame}
	
	% ---------------------------------------------------------------
	\begin{frame}[fragile]{Example: Old-Style Manual Configuration}
		\begin{lstlisting}[style=sparkstyle, caption={Manual Spark configuration in an old notebook}]
			from pyspark.sql import SparkSession
			
			spark = (SparkSession.builder
			.appName("LargeNotebookJob")
			.config("spark.executor.memory", "16g")
			.config("spark.executor.cores", "4")
			.config("spark.executor.instances", "10")
			.config("spark.driver.memory", "8g")
			.config("spark.sql.shuffle.partitions", "200")
			.getOrCreate())
		\end{lstlisting}
		
		\begin{itemize}
			\item Explicit values for memory, cores, and instances.
			\item Tuned for a specific large workload.
			\item This is what we might need to do --- or use a pre-configured notebook.
		\end{itemize}
	\end{frame}
	
	% ---------------------------------------------------------------
	\section{Self-Configured ``Large'' Notebooks}
	% ---------------------------------------------------------------
	
	\begin{frame}{Self-Configured Large Notebooks}
		\begin{block}{What are they?}
			Notebooks that come with \textbf{pre-set Spark configurations} optimized for large workloads.
		\end{block}
		
		\begin{itemize}
			\item The \textbf{tech team} creates different kinds of notebooks:
			\begin{itemize}
				\item Small / default notebooks for quick exploration.
				\item Medium notebooks for standard jobs.
				\item \textbf{Large notebooks} for heavy processing.
			\end{itemize}
			\item These notebooks are \textbf{self-configured} --- you do not need to manually set memory or executors.
			\item They are ready to use when you know your job is resource-intensive.
		\end{itemize}
		
		\begin{exampleblock}{Recommendation}
			Use a \textbf{large self-configured notebook} when your workload is known to be heavy.\\
			Fall back to manual configuration only if needed.
		\end{exampleblock}
	\end{frame}
	
	% ---------------------------------------------------------------
\begin{frame}[fragile]{How to Identify a Large Notebook}
	\begin{itemize}
		\item Check the notebook metadata or header comments.
		\item Look for pre-set Spark configuration cells:
	\end{itemize}
	
	\begin{lstlisting}[style=sparkstyle, caption={Self-configured large notebook (example)}]
		# This notebook is pre-configured for LARGE workloads.
		# Executor memory: 32g, Executors: 20, Cores: 8
		spark = SparkSession.builder.getOrCreate()
		# Configuration is applied automatically by the environment.
	\end{lstlisting}
	
	\begin{itemize}
		\item Ask the tech team which notebook profile to use.
		\item In many clusters, you can select the notebook type from a dropdown.
	\end{itemize}
\end{frame} 


	\section{Approach: Fix the Notebook}
	% ---------------------------------------------------------------
	
	\begin{frame}{Step-by-Step Plan}
		\begin{enumerate}
			\item \textbf{Reproduce the problem} --- run as-is and confirm the timeout.
			\item \textbf{Inspect old notebooks} --- find manual configs that worked.
			\item \textbf{Check available notebook profiles} --- ask tech team or look for large notebooks.
			\item \textbf{Switch to a large self-configured notebook} if available.
			\item \textbf{If not, manually set} executor memory, instances, cores, and driver memory.
			\item \textbf{Test with a smaller subset} first to validate the configuration.
			\item \textbf{Run the full workload} and monitor performance.
			\item \textbf{Iterate} --- hit and trial until stable and fast enough.
		\end{enumerate}
	\end{frame}
	
	% ---------------------------------------------------------------
\begin{frame}[fragile]{Manual Configuration: What to Change}
	
	\begin{lstlisting}[style=sparkstyle, caption={Manually tuning Spark for a large job}]
		spark = (SparkSession.builder
		.appName("FixedLargeNotebook")
		# Increase executor memory (e.g., 16g -> 32g)
		.config("spark.executor.memory", "32g")
		# Increase number of executors
		.config("spark.executor.instances", "20")
		# Increase cores per executor
		.config("spark.executor.cores", "8")
		# Increase driver memory if driver is collecting data
		.config("spark.driver.memory", "16g")
		# Tune shuffle partitions for large data
		.config("spark.sql.shuffle.partitions", "400")
		# Enable adaptive query execution (Spark 3+)
		.config("spark.sql.adaptive.enabled", "true")
		.getOrCreate())
	\end{lstlisting}
	
	\begin{itemize}
		\item Start with values from old notebooks that worked.
		\item Adjust based on cluster capacity and job behavior.
	\end{itemize}
	
\end{frame}
	% ---------------------------------------------------------------
	\begin{frame}{Hit and Trial: What to Monitor}
		\begin{itemize}
			\item \textbf{Spark UI} --- check stages, tasks, and spills.
			\item \textbf{Executor memory usage} --- are executors running out of memory?
			\item \textbf{Task duration} --- which stages are slow?
			\item \textbf{Shuffle read/write} --- excessive shuffle indicates partitioning issues.
			\item \textbf{GC time} --- high garbage collection means memory pressure.
		\end{itemize}
		
		\begin{block}{Iterative Tuning}
			\begin{itemize}
				\item Change one parameter at a time.
				\item Document what you changed and the result.
				\item Repeat until the job completes in acceptable time.
			\end{itemize}
		\end{block}
	\end{frame}
	
	% ---------------------------------------------------------------
	\section{Summary and Best Practices}
	% ---------------------------------------------------------------
	
	\begin{frame}{Summary}
		\begin{itemize}
			\item The notebook \textbf{times out} because default Spark resources are insufficient.
			\item Old notebooks solved this by \textbf{manually configuring} Spark.
			\item The tech team now provides \textbf{self-configured large notebooks}.
			\item \textbf{Preferred approach:} use a large notebook profile.
			\item \textbf{Fallback:} manually set executor memory, instances, cores, and driver memory.
			\item Expect \textbf{hit and trial} --- tuning is iterative.
		\end{itemize}
		
		\begin{exampleblock}{Key Takeaway}
			Match the Spark configuration to the workload.\\
			Use large self-configured notebooks when available; otherwise, tune manually.
		\end{exampleblock}
	\end{frame}
	
	% ---------------------------------------------------------------
	\begin{frame}{Best Practices}
		\begin{enumerate}
			\item \textbf{Know your workload} --- small, medium, or large.
			\item \textbf{Prefer self-configured notebooks} provided by the tech team.
			\item \textbf{Document} any manual configuration you apply.
			\item \textbf{Test with a subset} before running the full job.
			\item \textbf{Monitor Spark UI} during execution.
			\item \textbf{Iterate} --- do not expect perfect settings on the first try.
			\item \textbf{Communicate} with the tech team about cluster limits.
		\end{enumerate}
	\end{frame}
	
	% ---------------------------------------------------------------
	\begin{frame}{References and Further Reading}
		\begin{itemize}
			\item Apache Spark Documentation: \url{https://spark.apache.org/docs/latest/configuration.html}
			\item Spark Tuning Guide: \url{https://spark.apache.org/docs/latest/tuning.html}
			\item Internal tech team documentation on notebook profiles.
			\item Old notebooks in the repository (3--5 years old) --- examples of manual tuning.
		\end{itemize}
		
		\vspace{1em}
		\begin{center}
			\Large Questions?
		\end{center}
	\end{frame}
	
\end{document}